\documentclass{ip-lab}
\usepackage{booktabs}

\lablevel{n2}
\labnumber{Laboratorio 2}
\labtitle{Condicionales y cadenas de caracteres}

\makeheader

\begin{document}

\maketitle

\begin{sectionbox}{Preparación del ambiente de trabajo}
  Cree un nuevo módulo de funciones llamado \texttt{condicionales.py} en el que escribirá las funciones correspondientes a los puntos 1 a 2 de este laboratorio.
\end{sectionbox}

\begin{sectionbox}{Actividad 1: Validador de contraseñas}
Escriba una función que reciba como parámetro una contraseña y retorne un puntaje que indica qué tan segura es la contraseña según los siguientes criterios:

\begin{enumerate}
    \item Si la contraseña tiene al menos 8 caracteres de longitud, suma tantos puntos como caracteres.
    \item Si la contraseña contiene al menos una letra minúscula, suma 3 puntos.
    \item Si la contraseña contiene al menos una letra mayúscula, suma 5 puntos.
    \item Si la contraseña contiene al menos un dígito (0-9), suma 7 puntos.
\end{enumerate}

Debe contemplar que si se transgrede alguno de los siguientes criterios, la contraseña se considera insegura y el puntaje final es 0:
\begin{enumerate}
    \item No debe contener menos de 8 caracteres.
    \item No debe contener espacios en blanco.
    \item No debe consistir de solo dígitos.
    \item No debe ser igual a las siguientes contraseñas comunes: \texttt{password}, \texttt{admin}.
\end{enumerate}
\end{sectionbox}

\begin{sectionbox}{Actividad 2: Mejor contraseña}
Escriba una función que reciba como parámetros cuatro contraseñas y retorne un mensaje indicando cuál es la contraseña más segura de las cuatro y su valor.

Mida el puntaje de cada contraseña conforme a los criterios de la Actividad 1. 

En caso de empate, elija la primera contraseña con el puntaje más alto.
\end{sectionbox}

\pagebreak

\begin{sectionbox}{Actividad 3: Pruebas}
Verifique que su función de puntuación produce los puntajes de la siguiente tabla. Luego, use combinaciones de estas contraseñas para probar que compara contraseñas correctamente, incluyendo casos con puntajes iguales.

\begin{center}
  \begin{tabular}{lc}
    \toprule
    \textbf{Contraseña}  & \textbf{Puntaje} \\
    \midrule
    \texttt{1234567}     & 0                \\
    \texttt{12345678}    & 0                \\
    \texttt{mi clave}    & 0                \\
    \texttt{password}    & 0                \\
    \texttt{abcdefgh}    & 11               \\
    \texttt{PASSWORD}    & 13               \\
    \texttt{ABCDEFGH}    & 13               \\
    \texttt{Abcdefgh}    & 16               \\
    \texttt{1234567a}    & 18               \\
    \texttt{admin123}    & 18               \\
    \texttt{password1}   & 19               \\
    \bottomrule
  \end{tabular}
\end{center}
\end{sectionbox}

\begin{sectionbox}{Entrega}
Entregue el archivo a través de Bloque Neón en el laboratorio del Nivel 2 designado como ``N2-L2: Condicionales y cadenas de caracteres''.
\end{sectionbox}

\end{document}
